% 作者题证的编辑转录副本，不是独立下载的上游原件。
% 来源：https://github.com/BenoitDionne/PDE/blob/PDE/elliptic.tex

\paragraph{本节的环境与记号（据作者相邻段落整理）。}
$\Omega$为有界域，具有使$H^k(\Omega)\hookrightarrow L^2(\Omega)$紧的延拓/边界正则性；
$k\geq1$，$H^k_0(\Omega)\subseteq X\subseteq H^k(\Omega)$且$X$闭。
$\langle v,u\rangle_2=\int v\overline u$，$B:X\times X\to\mathbb C$为有界半双线性形式，线性于第一变量。
弱强制性为$\operatorname{Re}B(u,u)\geq C\|u\|_k^2-\lambda\|u\|_2^2$，$C,\lambda>0$。
\begin{theorem}[ell\_exist\_th2]
Let
\[
V=\{u\in X:B(v,u)=0\text{ for all }v\in X\},\qquad
W=\{w\in X:B(w,v)=0\text{ for all }v\in X\}.
\]
Then $\dim V=\dim W<\infty$. For $f\in L^2(\Omega)$, the variational
problem
\[
B(v,u)=\langle v,f\rangle_2\quad(v\in X),\qquad u\in X,
\]
has a solution if and only if $\langle w,f\rangle_2=0$ for all $w\in W$.
Solutions are unique up to addition of an element of $V$.
In particular, if $V=W=\{0\}$, each $f$ has a unique solution.
\end{theorem}
\paragraph{作者在主证明之前构造的解算子与伴随。}
If $B$ is coercive, Lax--Milgram provides a bounded linear operator
$F_B:L^2(\Omega)\to X$ satisfying
$B(v,F_Bf)=\langle v,f\rangle_2$. Here the functional is bounded because
$|\langle v,f\rangle_2|\leq\|f\|_2\|v\|_k$.
Define $B^*(u,v)=\overline{B(v,u)}$; it is likewise bounded and coercive,
so it has $F_{B^*}$. If $i:X\hookrightarrow L^2(\Omega)$ is the compact
inclusion, then $K_1=iF_B$ and $K_2=iF_{B^*}$ are compact, and
\begin{align*}
\langle K_1g,f\rangle_2
&=B^*(F_Bg,F_{B^*}f)
=\overline{B(F_{B^*}f,F_Bg)}\\
&=\overline{\langle F_{B^*}f,g\rangle_2}
=\langle g,K_2f\rangle_2.
\end{align*}
Thus $K_1^*=K_2$.
\begin{proof}[Proof of the theorem]
Choose $C>0,\lambda>0$ in the weak coercivity bound. Then
\[
\widetilde B(v,u)=B(v,u)+\lambda\langle v,u\rangle_2
\]
is coercive. By the preceding construction, now for $\widetilde B$,
there is a compact $K_1:L^2(\Omega)\to L^2(\Omega)$ with range in $X$
and
\[
\widetilde B(v,K_1f)=\langle v,f\rangle_2.
\]
For $u\in X$ we have the equivalences
\begin{align*}
 B(v,u)=\langle v,f\rangle_2\ (v\in X)
&\Longleftrightarrow\widetilde B(v,u)=\langle v,f+\lambda u\rangle_2\ (v\in X)\\
&\Longleftrightarrow u=K_1(f+\lambda u)\\
&\Longleftrightarrow \lambda^{-1}u-K_1u=\lambda^{-1}K_1f.
\end{align*}
Taking $f=0$ shows
\[
V=\ker(K_1-\lambda^{-1}I).
\]
Indeed an $L^2$ vector satisfying $u=\lambda K_1u$ automatically belongs
to $X$. The same construction for $B^*$ gives compact $K_2=K_1^*$ with
\[
W=\ker(K_2-\lambda^{-1}I).
\]
The compact-operator Fredholm theorem gives equal finite dimensions
of these two kernels, proving the first assertion.

The transformed equation is solvable exactly when
$K_1f\in\im(K_1-\lambda^{-1}I)$.
This range is closed by the same compact-operator theorem. The Hilbert
space range--kernel orthogonality identities give
\[
\im(K_1-\lambda^{-1}I)
=\ker(K_1^*-\lambda^{-1}I)^\perp=W^\perp.
\]
For $w\in W$,
\[
\langle K_1f,w\rangle_2
=\langle f,K_1^*w\rangle_2
=\lambda^{-1}\langle f,w\rangle_2.
\]
Therefore $K_1f\in W^\perp$ if and only if $f\in W^\perp$.
This is the solvability criterion. The difference of two solutions is
in $V$ by subtraction, and adding an element of $V$ preserves a solution.
The final assertion follows when the two kernels vanish.
\end{proof}
